\subsection{闭集}

定义 7. 在拓扑空间 X 中，X 的开集的余集称为闭集。

取余集，可知公理 \((\mathbf{O}_{\mathbf{I}})\) 与 \((\mathbf{O}_{\mathbf{II}})\) 取下列形式：

\((\mathbf{O}_{\mathbf{I}}^{\prime})\) 闭集的任何交是闭集。

\((\mathbf{O}_{\mathrm{II}}^{\prime})\) 闭集的任何有限并是闭集。

空集与整个空间 X 都是闭集（从而既开又闭；参见 § 11）。

在有理直线上，每个形如 \([a, \rightarrow[\) 的区间都是闭集，因为它的余集 \(]\leftarrow, a[\) 是开集；同样，每个形如 \(]\leftarrow, a]\) 的区间都是闭集；从而每个有界闭区间 \([a, b]\) 也是闭集，因为它是区间 \([a, \rightarrow[\) 与 \(]\leftarrow, b]\) 之交。

有理整数集 Z 在有理直线中是闭集，因为它的

余集 \(\bigcup_{n\in \mathbf{Z}}]n,n + 1[\) 是开集。

设 \((\mathbf{F}_i)_{i\in I}\) 是 \(\mathbf{A}\)（拓扑空间 \(\mathbf{X}\) 的一个子集）的覆盖；如果每个 \(\mathbf{F}_i\) 在 \(\mathbf{X}\) 中都是闭集，则称该覆盖为闭覆盖。

同胚 \(f\)（由拓扑空间 \(\mathbf{X}\) 到拓扑空间 \(\mathbf{X}'\) 上，第 1 目）可以刻画为 \(\mathbf{X}\) 到 \(\mathbf{X}'\) 上的一个双射，使得在 \(f\) 之下，\(\mathbf{X}\) 的每个闭子集的像是 \(\mathbf{X}'\) 的闭子集，并且在 \(f\) 之下，\(\mathbf{X}'\) 的每个闭子集的逆像是 \(\mathbf{X}\) 的闭子集。

